\documentclass[12pt]{article}
\usepackage[utf8]{inputenc}
\usepackage{amsmath, amssymb, amsthm}
\usepackage{geometry}

% Set margins to match a standard homework document
\geometry{a4paper, margin=1in}

\title{Recitation 3}
\parskip=1ex
\author{Math 311 Fall 2026 \\ Rutgers University}
\date{}

\begin{document}

\maketitle

\noindent{} \textbf{Question 1}

Let $A, B \subseteq (0, \infty)$ be non-empty and bounded above. Set
\[
    \alpha = \sup A, \qquad \beta = \sup B,
\]
and define
\[
    C =\{xy : x \in A, y \in B\}
\]
By definition of supremum,
\[
    \alpha \geq x : \forall x \in A
\]
and
\[
    \beta \geq y : \forall y \in B.
\]
Then, since all $x$ and $y$ are positive,
\[
    \alpha \beta \geq xy : \forall x \in A, \forall y \in B.
\]
That is, $\alpha \beta$ is an upper bound on $C$.

Suppose $c$ is an arbitrary upper bound on $C$.
Then
\[
    c \geq xy
\]
for every $x \in A$ and $y \in B$.
So
\[
    \frac{c}{y} \geq x.
\]
Then $\dfrac{c}{y}$ is an upper bound on $A$.
Since $\alpha = \sup A$,
\[
    \frac{c}{y} \geq \alpha,
\]
and
\[
    c \geq \alpha y.
\]
Then
\[
    \frac{c}{\alpha} \geq y
\]
That is, $\dfrac{c}{\alpha}$ is an upper bound on $B$.
Since $\beta = \sup B$,
\[
    \frac{c}{\alpha} \geq \beta,
\]
and
\[
    c \geq \alpha \beta.
\]
Therefore, $\alpha \beta$ is also the least upper bound on $C$ or
\[
    \alpha \beta = \sup C.
\]

\noindent{} \textbf{Question 2}
\begin{enumerate}
    \item [(a)] Let $S$ be an infinite set.
          We construct distinct elements
          \[
              a_1,a_2,a_3,\ldots\in S
          \]
          inductively.
          Since $S$ is infinite, it is also non-empty.
          Thus, we can select $a_1 \in S$ as the first element.

          Suppose that $a_1, a_2, ... , a_n$ have been chosen for $n \in \mathbb{N}$.
          Then
          \[
              \{a_1, a_2, ... , a_n\}
          \]
          is still a finite set. Since $S$ is infinite,
          \[
              S \setminus \{a_1, a_2, ... , a_n\} \neq \emptyset
          \]
          So we can choose $a_{n+1} \in S \setminus \{a_1, a_2, ... , a_n\}$.
          Therefore, by induction, we can choose countably infinitely many distinct elements
          \[
              A = \{a_1,a_2,a_3,\ldots\} \subseteq S.
          \]
          Define this function $f: \mathbb{N} \to A$ as
          \[
              f(n) = a_n.
          \]

          Suppose $f(m) = f(n)$. Then
          \[
              a_n = a_m.
          \]
          Since we define $a_i$ such that each element is distinct,
          it follows that $m = n$. Thus $f$ is injective.

          By definition of $a_i$, there is some $n \in \mathbb{N}$
          for each distinct element $a_i$. So $f$ is surjective.

          Therefore, $f$ is bijective, and
          \[
              \mathbb{N} \sim A \subseteq S.
          \]
    \item [(b)]
          Let $A = \{a_1,a_2,a_3,\ldots\} \subseteq S$ be the countably
          infinite set constructed in part (a).
          Define $F: S \to S \setminus \{a_1\}$ as
          \[
              F(x)=
              \begin{cases}
                  a_{n+1}, & x=a_n\text{ for some }n\in\mathbb N, \\
                  x,       & x\notin A.
              \end{cases}
          \]

          Suppose $F(x) = F(y)$.

          If $x,y \in A$, then
          \[
              x = a_m \qquad \text{and} \qquad y = a_n
          \]
          for some $m,n \in \mathbb{N}$.
          So $F(x) = F(y)$ implies that
          \[
              a_{m+1} = a_{n+1}.
          \]
          Since elements of $A$ are distinct, it follows that
          \[
              m = n.
          \]
          So
          \[
              x=y.
          \]

          If $x, y \notin A$, then $F(x) = x$ and $F(y) = y$.
          So $F(y) = F(x)$ implies $x = y$.

          What about if $x \in A$ and $y \notin A$, or the other way around?
          Suppose, for contradiction, that $F(x) = F(y)$.
          Since $x \in A $, it follows that $x = a_{i}$ for some $i \in \mathbb{N}$.
          Then $F(x) = a_{i+1}$, but $F(y) = y \notin A$.
          So $F(x) \neq F(y)$, contradicting our original assumption.
          Hence this case is impossible.

          Thus, $F$ is injective.

          Let
          \[
              y \in S \setminus \{a_1\}.
          \]
          If $y \in A$, then $y = a_n$ for some $n \geq 2$.
          Then
          \[
              F(a_{n-1}) = a_n = y.
          \]
          If $y \notin A$, then
          \[
              F(y) = y.
          \]
          Thus $F$ is surjective.

          Therefore $F$ is a bijection, so
          \[
              S \sim S \setminus \{a_1\}.
          \]
          Hence, every infinite set $S$ contains a proper subset that is equivalent to $S$.



    \item [(c)]
          Let $A \subseteq B$, where $A$ is countable and $B$ is uncountable.
          Suppose, for contradiction, that $B \setminus A$ is countable.
          Since $A$ is countable, then the union
          \[
              B = A \cup (B \setminus A)
          \]
          must be countable as well. But this contradicts the assumption that $B$ is uncountable.
          Hence $B \setminus A$ is uncountable.

          Since $B \setminus A$ is infinite, by part (a), we can choose distinct elements
          \[
              d_1, d_2, d_3, \ldots \in B \setminus A.
          \]
          Define $h: B \to B \setminus A$ as
          \[
              h(x)=
              \begin{cases}
                  d_{2n-1}, & x=a_n,                           \\
                  d_{2n},   & x=d_n,                           \\
                  x,        & x\notin A\cup\{d_1,d_2,\ldots\}.
              \end{cases}
          \]

          Suppose $h(x) = h(y)$.
          If $x, y \in A$, then
          \[
              d_{2n-1} = d_{2m-1}
          \]
          for some $m, n \in \mathbb{N}$.
          So $m = n$, and $x = y = a_n$.

          Similarly, if $x, y \in \{d_1, d_2, \ldots \}$, then
          \[
              d_{2n} = d_{2m}
          \]
          for some $m, n \in \mathbb{N}$.
          So $m = n$, and $x = y = d_n$.

          If $x, y \notin A \cup \{d_1, d_2, \ldots \}$, then
          \[
              y = x.
          \]

          If \(x\) and \(y\) belong to different cases, then \(h(x)\neq h(y)\) since each
          of the three cases have disjoint image sets.

          Thus, $h$ is injective.


          Let $y \in B \setminus A$.
          Then we have three possibilities.
          If $y = d_{2n-1}$, then choosing $x = a_n$, we have
          \[
              h(a_n) = d_{2n-1} = y.
          \]
          If $y = d_{2n}$, then choosing $x = d_n$, we have
          \[
              h(d_n) = d_{2n} = y.
          \]
          If $y \notin \{d_1, d_2, \ldots\}$, then choosing $x = y$, we have
          \[
              h(y) = y.
          \]
          Thus $h$ is surjective.

          Hence $h$ is a bijection, and so
          \[
              B \sim B \setminus A.
          \]

          If $A$ was instead finite, we simply shift each $d_i$ an extra $k$ elements, where $k$ is the length of $A$.

          e.g.
          \[
              h(x)=
              \begin{cases}
                  d_{i},   & x=a_i,                           \\
                  d_{n+k}, & x=d_n,                           \\
                  x,       & x\notin A\cup\{d_1,d_2,\ldots\}.
              \end{cases}
          \]
\end{enumerate}

\noindent{} \textbf{Question 3}
Let $\mathcal{F} : \{g: g : \mathbb{R} \to \mathbb{R}\}$
denote the set of all real-valued functions whose domain is $\mathbb{R}$.
Assume, for contradiction, that there exists a bijection
\[
    f: \mathbb{R} \to \mathcal{F}.
\]
For each $a \in \mathbb{R}$, let $g_a = f(a)$ be the function corresponding to the real number $a$.

Now define $h: \mathbb{R} \to \mathbb{R}$ by
\[
    h(x) = 1 + g_x(x).
\]
Notice that $h \in \mathcal{F}$.
Let $a \in \mathbb{R}$. Then
\[
    h(a) = 1 + g_a(a).
\]
Hence
\[
    h(a) \neq g_a(a),
\]
and since $a$ is arbitrary,
\[
    h \neq g_a
\]
for every $a \in \mathbb{R}$.

Since $g_a = f(a)$, this means that $h$ is never equal to $f(a)$, and so it is not in the image of $f$.
But $h \in \mathcal{F}$, so this contradicts the assumption that $f$ is surjective on $\mathcal{F}$.
Therefore, there is no bijection from $\mathbb{R}$ to $\mathcal{F}$, and so
\[
    \mathcal{F} \not \sim \mathbb{R}.
\]

\noindent{} \textbf{Question 4}
Let $f: (0,1) \to \mathbb{R}$.
Assume there is a constant $M >0$ such that for every finite set of distinct points
$x_1, ... , x_k \in (0, 1)$, we have
\[
    \sum_{j=1}^{k}|f(x_j)| \leq M.
\]
Let $S = \{x \in (0,1) : f(x) \neq 0\}$.
For $n \in \mathbb{N}$, define
\[
    S_n = \{ x \in (0,1): |f(x)| \geq \frac{1}{n} \}.
\]
Suppose, for contradiction, that some $S_n$ is infinite.
Then, we can choose infinitely many distinct elements from $S_n$, say
\[
    x_1, ... , x_k \in S_n
\]
for some $k \in \mathbb{N}$.
Since each $x_i \in S_n$,
\[
    |f(x_i)| \geq \frac{1}{n}.
\]
So
\[
    \sum_{i=1}^{k}|f(x_i)| \geq \sum_{i=1}^{k}\frac{1}{n} = \frac{k}{n}.
\]
By Archimedean Property, choose $k$ large enough so that
\[
    k > Mn,
\]
and so
\[
    \frac{k}{n} > M.
\]
Then
\[
    \sum_{i=1}^{k}|f(x_i)| > M,
\]
contradicting our assumption that
\[
    \sum_{j=1}^{k}|f(x_j)| \leq M.
\]
Thus, every $S_n$ must be finite.

Let
\[
    x \in \bigcup_{n=1}^{\infty} S_n.
\]
Then $x \in S_n$ for some $n \in \mathbb{N}$.
So $x \in (0, 1)$ and
\[
    |f(x)| \geq \frac{1}{n} > 0.
\]
Then $f(x) \neq 0$, so $x \in S$.
Hence
\[
    \bigcup_{n=1}^{\infty} S_n \subseteq S.
\]

Conversely, let
\[
    x \in S.
\]
So $x \in (0, 1)$ and $f(x) \neq 0$.
That is,
\[
    |f(x)| > 0.
\]
By Archimedean Property, there exists $n \in \mathbb{N}$ such that
\[
    \frac{1}{n} < |f(x)|.
\]
Then
\[
    |f(x)| \geq \frac{1}{n}.
\]
So $x \in S_n$ for some $n \in \mathbb{N}$.
Furthermore,
\[
    x \in \bigcup_{n=1}^{\infty} S_n.
\]
Hence
\[
    S \subseteq \bigcup_{n=1}^{\infty} S_n.
\]
Therefore
\[
    S = \bigcup_{n=1}^{\infty} S_n.
\]
Since $S_n$ is finite, then $S$ is the countable union of finite sets.
Hence $S$ must be countable as well.


\noindent{} \textbf{Question 5}
Let $C \subseteq [0,1]$ be uncountable, and define
\[
    E = \{a \in (0,1) : C \cap [a,1] \text{ is uncountable}\}.
\]
\begin{enumerate}
    \item [(a)]
          Suppose, for contradiction, that $E = \emptyset$.
          Then $C \cap [a, 1]$ is countable for every $a \in (0,1)$.
          Specifically, for each $n \geq 2$,
          \[
              C \cap [\frac{1}{n}, 1]
          \]
          is countable.
          Then
          \[
              C \cap (0, 1] = \bigcup_{n=2}^{\infty} (C \cap [\frac{1}{n}, 1]).
          \]
          Notice that the right side is the countable union of countable sets.
          So $C \cap (0, 1]$ is countable.
          Adding the finite element $0$ means $C$ is countable as well.
          However, this contradicts our assumption that $C$ is uncountable.

          Therefore, $E$ is non-empty.

    \item [(b)]
          Let $\alpha = \sup E$.
          Consider
          \[
              C = [0, 1).
          \]
          Then $C$ is indeed an uncountable subset of $[0,1]$.
          Then, for every $a \in (0, 1)$
          \[
              C \cap [a, 1] = [a, 1).
          \]
          Which is also uncountable.
          Thus, $E = (0, 1)$, and $\alpha = 1$.
          But
          \[
              C \cap [\alpha, 1] = \emptyset,
          \]
          which is clearly not uncountable.

          Therefore, $C \cap [\alpha, 1]$ does not have to be uncountable.

    \item [(c)]
          False, take
          \[
              C = \left\{ \frac{1}{n} : n \in \mathbb{N} \right\}.
          \]
          This set is infinite.
          For any $a \in (0, 1)$, if $\frac{1}{n} \in [a, 1]$, then $\frac{1}{n} \geq a$.
          So $n \leq \frac{1}{a}$.
          Then only a finite amount of $n \in \mathbb{N}$ will meet this requirement.
          Thus $C \cap [a, 1]$ is finite for every $a \in (0,1)$.
          Therefore, $E = \emptyset$.
\end{enumerate}

\end{document}

